% Auto-generated by tools/qbe.py project-article-update.
% Do not edit this file by hand; edit the proof artifacts or article sections instead.

\section{Generated Current Cycle Status}
\label{app:generated-cycle-status}

This appendix page is generated from the latest ABEIS proof cycle.  It is a
status bridge for article maintenance, not a polished theorem proof.  Stable
mathematical claims should be copied into the main text only after the
corresponding Lean declaration, source-paper anchor, and proof-note export are
available.

\paragraph{Cycle.}
Task \texttt{QBE-AUTO-002}, cycle \texttt{1}, run
\texttt{runs/20260615-022953-QBE-AUTO-002-cycle01}.  Generated at
\texttt{2026-06-15 02:41:19}.

\paragraph{Lean status signal.}
\begin{itemize}
  \item \texttt{QuantumBlockEncoding/RobinMatrix.lean:25243:  sorry}
  \item \texttt{QuantumBlockEncoding/RobinMatrix.lean:25273:  sorry}
\end{itemize}

\paragraph{Plain-language status.}
The current GHL case study is not blocked because the paper lacks a proof sketch.  It is blocked because ABEIS must turn the circuit in the paper into an exact matrix statement and prove that the clean block entry has the coefficient claimed by the paper.  In ordinary terms, the paper says that a specific path through the circuit leaves the desired coefficient, while Lean requires the corresponding matrix entry and all unwanted branches to be proved exactly.

\paragraph{Pre-Lean verifier candidates.}
These checks are necessary-condition filters, not proofs.  A failed exact
finite check usually indicates that the current Lean target, circuit
transcript, or index map is wrong.  A passing check only means the candidate
survived a cheaper diagnostic; the final claim still requires a named Lean
theorem or an explicit contract.
\begin{itemize}
  \item \textbf{active [0,0] entry.} Fast check: exact rational matrix or path-sum evaluation.  Necessary because if the exact finite entry is not the target coefficient, the Lean equality for that entry cannot be true.  Lean still must prove: a passing check is only a counterexample filter; Lean must still prove the named entry lemma.
  \item \textbf{remaining branch vanish/cancel.} Fast check: support and path checker for the remaining backend slots.  Necessary because if an unwanted clean-branch path survives numerically, the block projection cannot match the paper target.  Lean still must prove: Lean must still prove the zero/cancellation lemma in the formal circuit semantics.
  \item \textbf{Ry boundary convention.} Fast check: symbolic 2-by-2 rotation check.  Necessary because a mismatched half-angle convention changes the boundary amplitude before any Lean tactic is relevant.  Lean still must prove: Lean must still record the convention bridge as a source-supported theorem.
  \item \textbf{sparse-access map.} Fast check: finite range/injectivity/permutation check.  Necessary because a reversible oracle cannot exist for a colliding or out-of-range finite map.  Lean still must prove: Lean must still prove the reversible extension and cleanup obligations.
  \item \textbf{coefficient oracle clean branch.} Fast check: exact finite evaluator for f(x\_i)/N\_f.  Necessary because the final block entry uses this coefficient, so a wrong clean branch invalidates the target theorem.  Lean still must prove: Lean must still prove bounds, orthogonality, and unitary completion or keep them as contracts.
\end{itemize}

\paragraph{Current proof-DAG frontier.}
\begin{itemize}
  \item 1. lower1 writes a source-dependency/proof-DAG correction packet for the projection/backend-expansion frontier. It must explain whether the paper supports a corrected source-facing branch-sum proposition, or whether this is a source-contract gap.
  \item 2. lower3 records a necessary-condition verifier packet for the raw target: the existing all-one selected-branch counterexample means the unchanged \textasciigrave{}backendExpansionStatement\textasciigrave{} is not a valid lower2 theorem target.
  \item 3. lower2 does not try to prove the refuted proposition unchanged. If lower2 is run before lower1/lower3 finish the correction, it should make no Lean edit and log \textasciigrave{}error\_class=finite\_matrix\_counterexample\textasciigrave{}. After the correction packet exists, lower2 may ed...
\end{itemize}

\paragraph{Open obligation signal.}
\begin{itemize}
  \item source-prepared product/projection packet: Lean \textasciigrave{}oneTermRobinGamma3BoundarySourcePreparedProductProjectionObligation\_n3\textasciigrave{}; transcript theorem; class QBE-local proof-DAG bookkeeping leaf; status compiled; stale as lower2 target
  \item source-prepared projection to backend fold: Lean \textasciigrave{}oneTermRobinGamma3BoundarySourcePreparedProjection\_to\_backendFold\_n3\textasciigrave{}; class wrapper over compiled source-prepared backend evaluator under \textasciigrave{}hUniform\textasciigrave{}; status compiled; stale as lower2 target
  \item backend fold to slot-\textasciigrave{}2\textasciigrave{} projected product: Lean \textasciigrave{}oneTermRobinGamma3BoundaryBackendFold\_to\_slot2ProjectedProduct\_n3\textasciigrave{}; class backend fold collapse plus selected-slot product evaluator under \textasciigrave{}hentry\textasciigrave{}; status compiled; stale as lower2 target
  \item source-prepared projection to slot-\textasciigrave{}2\textasciigrave{} projected product: Lean \textasciigrave{}oneTermRobinGamma3BoundarySourcePreparedProjection\_slot2\_to\_projectedBranchProduct\_n3\textasciigrave{}; class composition of the previous two bridge leaves; status compiled; stale as lower2 target
  \item fixed product-to-coefficient theorem for \textasciigrave{}(0,0)\textasciigrave{}: Lean \textasciigrave{}oneTermRobinGamma3ProductToCoefficientObligation 3 0 0\textasciigrave{}; class coefficient equality plus normalizer algebra and boundary coefficient convention; status open; blocked on corrected projection/backend branch-sum route
  \item unchanged backend-expansion raw target: Lean \textasciigrave{}oneTermRobinGamma3BoundaryBackendBranchContributionTarget\_n3.backendExpansionStatement\textasciigrave{}; class contract drift / finite counterexample for current theorem-facing route; status refuted; \textasciigrave{}oneTermRobinGamma3BoundaryBackendExpansionStatement\_not\_n3\textasciigrave{} compiled
  \item corrected projection/backend branch-sum leaf: Lean to be named by lower1/lower3 correction packet; class internal GHL step plus QBE-local finite projection lemma, or source-contract gap; status active correction frontier
  \item diagnostic raw equality route: Lean \textasciigrave{}oneTermRobinGamma3BoundaryUnitaryEntry\_eq\_backendFold\_n3\textasciigrave{}; \textasciigrave{}oneTermRobinGamma3BoundaryEvalGateMatrices\_eq\_sevenGateMatrix\_n3\textasciigrave{}; class existing diagnostic \textasciigrave{}sorry\textasciigrave{} route; status forbidden as dependency
\end{itemize}

\paragraph{Open GHL contribution obligations.}
\begin{itemize}
  \item \texttt{Uindic}: Indicator unitary; status Permutation and self-inverse helpers compile, but this is not yet the final block-encoding proof..
  \item \texttt{RyBoundary}: Boundary-controlled Ry rotations; status Active convention audit: the Ry angle convention must be source-supported before closing the boundary amplitude proof..
  \item \texttt{RobinTheorem}: One-term Robin block-encoding theorem; status Main active target: the theorem-facing Lean statement is not closed..
  \item \texttt{GammaSlices}: Gamma wavefunction slices and coefficient product; status Several finite boundary lemmas exist; the final projection/product bridge is still open..
  \item \texttt{FigRobin}: Figure 4 Robin circuit transcript; status Transcript guard compiles for visible indicator dagger and H preparation; the active seven-gate backend remains a component..
  \item \texttt{OneD}: One-dimensional Hamiltonian block encoding; status Deferred until the one-term Robin theorem closes..
  \item \texttt{MultiD}: Multidimensional block encoding; status Planned; depends on one-term and LCU abstractions..
\end{itemize}

\paragraph{Open external technical lemma obligations.}
\begin{itemize}
  \item \texttt{tl-ghl-lemma1-banded-sparse-access}: contract-only; next action Keep as explicit external contract until the exact cited construction is formalized or imported as a theorem..
  \item \texttt{tl-ghl-lemma3-sparse-amplitude}: contract-only; next action Maintain the clean-branch contract and isolate any missing unitarity proof as an external technical lemma..
  \item \texttt{tl-ghl-theorem5-piecewise-polynomial-of}: contract-only; next action Use only as a named contract in the current theorem; do not invent stronger smoothness or range assumptions..
  \item \texttt{tl-uniform-sparse-register-preparation}: contract-only; next action Keep as a theorem-facing contract unless the exact state-preparation proof is needed to close a resource theorem..
  \item \texttt{tl-ry-boundary-amplitude-convention}: obligation; next action Do not silently change the paper angle; prove the convention bridge or record the exact cited convention..
  \item \texttt{tl-qsvt-blockencoding-simulation}: paper-cited; next action Defer until the one-term Robin theorem and LCU combination are closed..
  \item \texttt{tl-clean-block-definition}: contract-only; next action Use as the final target predicate; ensure circuit entry lemmas are bridged to this semantic definition..
\end{itemize}

\paragraph{Recent typed verifier feedback.}
\begin{itemize}
  \item Leaf \texttt{backend\_expansion\_correction}, class \texttt{finite\_matrix\_counterexample}; next route: lower1/lower3 must refine the projection/backend target before lower2 Lean proof search; lower2 should not prove the raw backendExpansionStatement unchanged..
  \item Leaf \texttt{source\_projection\_slot2\_product}, class \texttt{symbolic\_bridge\_gap}; next route: advance to product\_to\_coefficient\_3\_0\_0 / oneTermRobinGamma3ProductToCoefficientObligation 3 0 0.
  \item Leaf \texttt{source\_projection\_slot2\_product}, class \texttt{symbolic\_bridge\_gap}; next route: Lower2 should add exactly one theorem, oneTermRobinGamma3BoundarySourcePreparedProjection\_slot2\_to\_projectedBranchProduct\_n3, by calc-transitivity through oneTermRobinGamma3BoundarySourcePreparedProjection\_to\_backendFold\_n3 H env hUniform and oneTermRobinGamma3BoundaryBackendFold\_to\_slot2ProjectedProduct\_n3 env hentry..
  \item Leaf \texttt{source\_projection\_slot2\_product}, class \texttt{symbolic\_bridge\_gap}; next route: Prove oneTermRobinGamma3BoundarySourcePreparedProjection\_slot2\_to\_projectedBranchProduct\_n3 by composing oneTermRobinGamma3BoundarySourcePreparedProjection\_to\_backendFold\_n3 H env hUniform with oneTermRobinGamma3BoundaryBackendFold\_to\_slot2ProjectedProduct\_n3 env hentry..
  \item Leaf \texttt{backend\_fold\_to\_slot2\_projected\_product}, class \texttt{None}; next route: compose oneTermRobinGamma3BoundarySourcePreparedProjection\_slot2\_to\_projectedBranchProduct\_n3 from the two compiled bridge leaves.
  \item Leaf \texttt{backend\_fold\_to\_slot2\_projected\_product}, class \texttt{symbolic\_bridge\_gap}; next route: Lower2 should prove exactly one theorem: oneTermRobinGamma3BoundaryBackendFold\_to\_slot2ProjectedProduct\_n3 by calc-transitivity through oneTermRobinGamma3BoundaryBackendBranchFoldEval\_eq\_selectedSlotContribution\_n3 env and oneTermRobinGamma3BoundaryProjectionSummationObstruction\_selectedSlotEval\_n3 env hentry. After that, compose it with oneTermRobinGamma3BoundarySourcePreparedProjection\_to\_backendFold\_n3..
\end{itemize}

\paragraph{Recent proof-attempt memory.}
\begin{itemize}
  \item \texttt{proof-attempts/QBE-AUTO-002/backend-expansion-correction-middle-packet-20260615-0233.md}
  \item \texttt{proof-attempts/QBE-AUTO-002/source-prepared-projection-product-composite-lower1-handoff-20260615-021428.md}
  \item \texttt{proof-attempts/QBE-AUTO-002/source-prepared-projection-product-composite-middle-packet-20260615-021428.md}
  \item \texttt{proof-attempts/QBE-AUTO-002/source-prepared-projection-product-composite-lower1-dag-20260615-015440.md}
  \item \texttt{proof-attempts/QBE-AUTO-002/source-prepared-projection-product-bridge-middle-packet-20260615-015440.md}
\end{itemize}

\paragraph{Article-writing rule.}
The project report may discuss this cycle as process evidence.  It must not
claim completion of the first case study \citep{guseynovHuangLiu2025} unless
the final theorem-facing block-extraction statement is closed in Lean and the
Markdown/LaTeX proof map has been synchronized.
